% !TeX root = ../../../main.tex
\section{Ask the Research Question Before Drawing the Circuit}\label{sec:contract}

\subsection{The bounded problem specification}

A research question should be written as a contract before it is written as a circuit. One useful specification is
\begin{equation}
\mathcal{P}=\bigl(X,\mathcal{D},f,C,S,\epsilon,\delta,\Lambda,B_C\bigr),
\label{eq:problem-tuple}
\end{equation}
where $X$ is the candidate or state space; $\mathcal{D}$ describes how inputs are accessed; $f$ is a predicate, score, likelihood, or operator; $C$ is the constraint set; $S$ records exploitable structure such as sparsity, symmetry, periodicity, locality, low rank, or topology; $\epsilon$ is the permitted numerical error; $\delta$ is the permitted failure probability; $\Lambda$ is the physical resource envelope; and $B_C$ is the strongest relevant classical baseline.

The corresponding quantum interface can be written
\begin{equation}
\mathcal{Q}=\bigl(A,O_f,U,M\bigr),
\label{eq:quantum-interface}
\end{equation}
where $A$ prepares the input state, $O_f$ implements the problem-dependent oracle or block encoding, $U$ produces the useful interference or dynamics, and $M$ defines the measured observable or sample. An abstract execution is therefore
\begin{equation}
\ket{0}\xrightarrow{A}\ket{\psi_{\mathrm{in}}}
\xrightarrow{O_f}\ket{\psi_f}
\xrightarrow{U}\ket{\psi_{\mathrm{out}}}
\xrightarrow{M} y.
\label{eq:pipeline}
\end{equation}

The tuples in \cref{eq:problem-tuple,eq:quantum-interface} are \emph{structured definitions}, not formal objects from which results are derived; their role is to force every assumption of an application claim onto the page. They prevent a common research failure: proving that $U$ is fast while quietly assuming that $A$ and $O_f$ are free, $M$ returns more information than measurement permits, and $B_C$ is a classroom algorithm from twenty years ago. A theorem can be correct under those assumptions. An application claim cannot hide them.

\subsection{Algorithm 1: the Quantum Opportunity Audit}\label{sec:algorithm}

Written as an executable procedure, the framework takes a candidate application and returns one of four decisions: \textsc{proceed}, \textsc{reframe}, \textsc{classical-only}, or \textsc{insufficient-evidence}. No numerical ``opportunity score'' is defined, deliberately: a gated sequence of falsifiable tests is more defensible than a weighted sum of judgments.

\begin{thesisbox}[Algorithm 1: Quantum Opportunity Audit]
\textbf{Inputs:} scientific question; target observable; tolerances $(\epsilon,\delta)$; deadline; hardware envelope $\Lambda_{\mathrm{HW}}$; best known classical workflow $B_C$.

\textbf{Gate 1 --- Contract.} Instantiate every field of $\mathcal{P}$ \eqref{eq:problem-tuple}. \emph{Stop condition:} any field cannot be written down $\Rightarrow$ \textsc{insufficient-evidence}.

\textbf{Gate 2 --- Mechanism.} Identify which mechanism in \cref{sec:mechanisms} the structure $S$ feeds, and the evidence class behind it (theorem, engineering approximation, or hypothesis). \emph{Stop condition:} no named mechanism $\Rightarrow$ \textsc{classical-only}.

\textbf{Gate 3 --- Interface.} Cost $A$, $O_f$, $M$ in \cref{eq:quantum-interface}, including uncomputation. \emph{Stop condition:} interface costs consume the claimed advantage $\Rightarrow$ \textsc{reframe} (change output contract or access model) or \textsc{classical-only}.

\textbf{Gate 4 --- Physics.} Run \cref{sec:resources}: error budget \eqref{eq:error-budget}, failure budget \eqref{eq:delta-union}, shot budget \eqref{eq:worstcase-shots} or \eqref{eq:neyman-shots}, noise triage \eqref{eq:nisq-screen} or fault-tolerant floor \eqref{eq:distance}. \emph{Stop condition:} no parameter regime fits $\Lambda_{\mathrm{HW}}$ and the deadline $\Rightarrow$ \textsc{classical-only} (revisit when hardware moves).

\textbf{Gate 5 --- Opponent.} Apply the baseline protocol below and estimate the crossover \eqref{eq:crossover}. \emph{Stop condition:} no projected crossover inside the envelope $\Rightarrow$ \textsc{classical-only}.

\textbf{Output:} the decision, the hybrid partition $\Pi^{*}$ (\cref{sec:hybrid}), and a predeclared evidence plan (\cref{sec:evidence}). \textsc{proceed} requires all five gates passed; the evidence plan then becomes binding.
\end{thesisbox}

Appendix~\ref{app:demo} executes this procedure end-to-end on a deliberately small, exactly solvable instance---where it correctly returns \textsc{classical-only}---and on the phase-encoding proposal of \cref{sec:audits}, which fails at Gate~3.

\paragraph{The gates are judgments, not measurements.} Each gate verdict is an auditable engineering judgment, not an observer-independent measurement. The procedure therefore requires, for every gate, a written record containing the evidence considered, the assumptions made, the identity of the classical baseline, the calculation performed, and the rationale for the decision. If qualified evaluators disagree on a gate and the disagreement cannot be resolved by that record, the audit returns \textsc{insufficient-evidence}; disagreement is never averaged into a favorable score. Inter-rater agreement for this procedure has not been measured: what the framework claims is procedural traceability---every verdict can be inspected and challenged---not evaluator invariance.

\begin{figure}[htbp]
\centering
\begin{tikzpicture}[
  node distance=7mm,
  every node/.style={font=\small},
  box/.style={rounded corners=1.5mm,draw=quantum,fill=mist,minimum width=0.84\textwidth,minimum height=9mm,align=center},
  arrow/.style={-{Latex[length=2.2mm]},thick,draw=quantumdark}
]
\node[box] (contract) {\textbf{Scientific contract:} observable, scale, $\epsilon$, $\delta$, deadline, and validation};
\node[box,below=of contract] (structure) {\textbf{Mechanism audit:} symmetry, spectrum, rare event, path interference, or native quantum dynamics};
\node[box,below=of structure] (interface) {\textbf{Interface audit:} state preparation, reversible oracle, compact output, and uncomputation};
\node[box,below=of interface] (budget) {\textbf{Physical audit:} topology, depth, errors, shots, error correction, classical latency, and baseline crossover};
\draw[arrow] (contract) -- (structure);
\draw[arrow] (structure) -- (interface);
\draw[arrow] (interface) -- (budget);
\end{tikzpicture}
\caption{The order of operations for a credible quantum research program. A circuit is designed only after the scientific and interface contracts are explicit.}
\label{fig:order}
\end{figure}

\subsection{Define the output before the state}

The output contract is the fastest way to expose a weak proposal. Quantum algorithms are naturally good at returning a sample, a phase, an expectation value, a low-dimensional statistic, a property, or a candidate that can be verified classically. They are generally poor at returning every entry of a large vector. The Harrow--Hassidim--Lloyd algorithm is the canonical warning: it prepares a state proportional to the solution of a structured linear system and efficiently estimates appropriate observables; it does not print an $N$-component solution in polylogarithmic time \citep{harrow2009}. Reading $N$ classical numbers still costs at least order $N$.

Before selecting an encoding, the researcher should answer four questions.

\begin{enumerate}
  \item What single scientific decision will be made from the output?
  \item Is the desired output a scalar, sample, property, witness, ranked list, or full object?
  \item How will correctness be checked when exact classical simulation is no longer available?
  \item What level of approximation changes the scientific conclusion rather than merely the last decimal place?
\end{enumerate}

The last question matters because precision is a resource. An algorithm with complexity $O(1/\epsilon)$ can be theoretically superior to a classical $O(1/\epsilon^2)$ Monte Carlo method, but not necessarily at the $\epsilon$ required by the application once constants, state preparation, and fault tolerance are included. ``Faster in precision'' is not the same as ``faster by Friday.''

\subsection{Define the classical opponent honestly}

The baseline is not brute force unless brute force is what competent practitioners actually use. A quantum optimization method should meet mixed-integer programming, constraint programming, local search, decomposition, learned heuristics, and GPU implementations on the same instance distribution. A quantum simulation proposal should meet tensor networks, Monte Carlo, perturbation theory, reduced models, and high-performance computing in the regime where each is strongest. A quantum machine-learning claim should meet compressed, randomized, kernel, and quantum-inspired classical methods. Tang's dequantization of recommendation-system speedups remains an important engineering lesson: strong quantum input assumptions may grant a classical competitor more power as well \citep{tang2019,tang2021}.

The comparison must fix the same problem size $N$, accuracy $\epsilon$, confidence $1-\delta$, hardware budget, preprocessing, and acceptable output. Only then does the desired inequality have meaning:
\begin{equation}
T_Q(N,\epsilon,\delta;\Lambda)+T_{\mathrm{interface}}
<T_C(N,\epsilon,\delta;B_C).
\label{eq:crossover}
\end{equation}

\paragraph{Baseline protocol.} ``Strongest relevant classical baseline'' becomes operational only when the comparison is specified in advance:
\begin{enumerate}
  \item name the classical implementations (solver, version, and configuration) that practitioners currently use for this problem class, including at least one state-of-the-art heuristic and one exact or certified method where feasible;
  \item grant both sides identical preprocessing, instance distributions, and hardware budgets (including graphics processing units for the classical side);
  \item give the classical side a tuning budget equal to the quantum side's calibration and transpilation effort;
  \item compare at the same $(\epsilon,\delta)$ and the same output contract, reporting time-to-solution distributions rather than best cases;
  \item state which costs are amortized on each side and over how many uses.
\end{enumerate}
